\documentclass[11pt]{article}
\usepackage[margin=1.2in]{geometry}
\usepackage{amsmath,amssymb,amsthm}
\newtheorem{theorem}{Theorem}
\newtheorem{corollary}[theorem]{Corollary}
\newtheorem{remark}[theorem]{Remark}
\newcommand{\FF}{\mathbb F}
\title{A reflection group cannot orient itself:\\
an elementary obstruction to deriving fermion chirality\\ from Weyl-group substrates}
\author{W.~Dahn\thanks{With machine-verified witnesses; see the repository
cited in \S\ref{sec:witness}. The constructions in \S\ref{sec:case} are due to
the repository's GAP track; the observation connecting them is the present
note.}}
\date{July 2026}
\begin{document}
\maketitle

\begin{abstract}
\noindent We record an elementary but easily-missed obstruction: if a physical
structure is built from a finite reflection group $W$, then no $W$-invariant
datum can select one of two chiralities exchanged by an element of determinant
$-1$ --- and a reflection group always contains such elements. Consequently any
programme deriving the standard model's handedness from a Weyl-group substrate
must import the choice from outside. We state the obstruction, give the
three-line proof, describe the worked case ($W(E_6)$ acting on the half-spin
pair of $D_5$, where every ingredient was machine-verified), and generalize:
every multiplicity such a substrate offers on which its symmetry acts
transitively is a torsor, hence internally unselectable. The obstruction is
cheap to check and, in the authors' experience, expensive not to.
\end{abstract}

\section{The obstruction}

\begin{theorem}\label{thm:main}
Let $W$ be a group acting on a structure $S$, and let $c_+,c_-$ be two objects
built from $S$ (``chiralities'') exchanged by some $w\in W$. Then no
$W$-invariant quantity distinguishes $c_+$ from $c_-$.
\end{theorem}

\begin{proof}
An invariant $F$ satisfies $F(c_-)=F(w\cdot c_+)=F(c_+)$.
\end{proof}

\begin{corollary}\label{cor:refl}
If $W$ is generated by reflections and the two chiralities are the half-spin
representations $S^\pm$ of an orthogonal structure on which $W$ acts with
$\det$ surjecting onto $\{\pm1\}$, the hypothesis of
Theorem~\ref{thm:main} holds automatically: improper elements exchange
$S^+$ and $S^-$ (the standard Pin/Spin fact), and a reflection group contains
improper elements --- its generators. Selecting a chirality is selecting a
coset of $\ker(\det)$, i.e.\ \emph{orienting} the root system, and a group
containing its own orientation reversals cannot do so.
\end{corollary}

The content is not the proof; it is noticing that the hypothesis applies, and
what that does to a research programme (\S\ref{sec:cost}).

\section{The worked case}\label{sec:case}

The case that prompted this note. A programme built on the symplectic
generalized quadrangle $W(3,3)$ --- whose automorphism group is
$W(E_6)\cong\mathrm{PGSp}(4,3)$, of order $51840$ --- constructs, in
characteristic zero, the split orthogonal module $V=5a\oplus5a^*$ over
$\mathbb Q(\omega)$ for $U_4(2)\cong\mathrm{PSp}(4,3)$ and its two
$16$-dimensional half-spin representations
$S^\pm=\Lambda^{\mathrm{even/odd}}(5a)$: genuinely complex, mutually
conjugate, non-isomorphic. The chirality is \emph{present}. The machine
verification then supplies:
\begin{enumerate}
\item an integral involution $T$ with $\langle U_4(2),T\rangle=W(E_6)$ and
$\det T=-1$;
\item three independent confirmations that $T$ exchanges $S^+$ with $S^-$:
the induced outer automorphism swaps the conjugate pair $5a\leftrightarrow
5a^*$; coefficient conjugation exchanges the two families of maximal
isotropics; and $\det T=-1$ is improper.
\end{enumerate}
By Theorem~\ref{thm:main}, no $\mathrm{PGSp}(4,3)$-invariant --- hence no
datum constructed from the quadrangle --- selects a half-spin. The substrate
\emph{hosts} the standard model's chirality and provably cannot \emph{choose}
it.

\begin{remark}
The deciding number, $\det T=-1$, had been computed and certified some fifteen
working iterations before anyone multiplied it by the Pin/Spin fact. The
obstruction is easy; \emph{noticing that it applies} is the entire difficulty,
which is why we recommend the check of \S\ref{sec:check}.
\end{remark}

\section{The generalization: multiplicities are torsors}

The same argument applies beyond chirality. In the worked case, every
multiplicity the substrate offers on which its symmetry acts transitively ---
two half-spins; three integral lattice classes cyclically permuted by the
Eisenstein unit $\omega$; twenty-seven machine states under a regular
Heisenberg action --- is a \emph{torsor}: a set with no distinguished point.
An invariant cannot distinguish points of a transitive orbit, so none of these
choices is internally selectable. A torsor is a group that has forgotten its
identity element; the substrate specifies everything except the base points,
and the missing physical input is exactly a \emph{section}, one per torsor.

Two refinements from the worked case sharpen the picture. First, the torsor
property can fail informatively: the twelve points collinear with a base point
of $W(3,3)$ admit \emph{no} regular group action at all (every involution in
the point stabilizer fixes one of them, by eigenspace dimension counts valid
for all odd $q$; the set has even size; Cauchy's theorem finishes it). The
substrate's multiplicities split into free torsors and blocked fibrations.
Second, the torsor identifications can be strong without being total: the
twenty-seven machine states match the twenty-seven lines of a cubic surface at
the level of the qutrit Pauli \emph{and} Clifford groups, while their invariant
geometries provably differ --- by exactly the central phase fiber.

\section{The cheap check}\label{sec:check}

For any programme deriving handedness (or any binary/n-ary vacuum choice) from
a structure with symmetry group $W$:
\begin{center}
\emph{Find the element of $W$ that exchanges your candidates.\\
If it exists, stop searching for a selector, and compute $\det$.}
\end{center}
If it exists, no internal quantity will ever select, and a search that feels
convergent is not. If the action on your candidates is transitive, they form a
torsor and the missing input is a section --- which at least tells you the
\emph{type} of the physics still owed.

\section{Witnesses}\label{sec:witness}

Every computational claim above is certified by machine-verified witnesses
(Python/GAP, with JSON certificates and a self-verifying claims ledger) in the
public repository \texttt{wilcompute/W33-Theory}, release
\texttt{v1.0-selection-layer-closed}; see in particular the witnesses numbered
346, 348, 354, 370--372 and the capstone note \emph{The 27-fold way of
$W(3,3)$}.

\end{document}
